% =============================================================
% บทที่ 4 - ผลการดำเนินงาน
% =============================================================

ในการทดลองประมวลผลและทำนายผลการแข่งขันฟุตบอลพรีเมียร์ลีกอังกฤษ ผู้จัดทำได้นำชุดข้อมูลการแข่งขันในฤดูกาล 2023/2024 จำนวนทั้งหมด 380 แมตช์ มาทำการทดสอบและวัดประสิทธิภาพของแบบจำลอง Machine Learning ที่ได้พัฒนาขึ้น

% -----------------------------------------------------------
\section{การทดสอบการทำงานของกระบวนการข้อมูล (Pipeline Validation)}
% -----------------------------------------------------------

การประมวลผลข้อมูลผ่าน Data Pipeline สามารถดำเนินงานได้อย่างสมบูรณ์ โดยไม่พบปัญหา Data Leakage จากการตรวจสอบลำดับเวลา ผลการจัดทำ Feature Extraction และการจัดการค่าที่ขาดหาย (Missing Values) ในชุดข้อมูลทดสอบแสดงสถานะตามตารางที่ \ref{tab:pipeline_status}

\begin{table}[H]
\centering
\caption{ผลการตรวจสอบความถูกต้องของระบบการประมวลผลข้อมูล}
\label{tab:pipeline_status}
\begin{tabular}{|l|c|l|}
\hline
\textbf{ขั้นตอนการทำงาน} & \textbf{สถานะ} & \textbf{รายละเอียด} \\ \hline
Entity Resolution (Name Mapping) & ผ่าน & จับคู่ชื่อทีมจาก 4 แหล่งข้อมูลสำเร็จ 100\% \\ \hline
Elo Rating Calculation & ผ่าน & คำนวณคะแนน Elo ต่อเนื่องตั้งแต่นัดแรกถึงนัดสุดท้าย \\ \hline
Rolling Feature Generation & ผ่าน & คำนวณค่าเฉลี่ย 5 และ 10 นัดย้อนหลังสำเร็จ \\ \hline
Temporal Data Splitting & ผ่าน & แบ่งข้อมูลตามวันแข่งจริง ไม่พบ Index ซ้ำซ้อน \\ \hline
\end{tabular}
\end{table}

% -----------------------------------------------------------
\section{ผลการจำแนกประเภทผลการแข่งขัน (Classification Results)}
% -----------------------------------------------------------

จากการนำแบบจำลอง Logistic Regression, Random Forest, XGBoost และ LightGBM มาทดลองทำนายผลการแข่งขันจำนวน 380 แมตช์ ในชุดข้อมูล Test Set ได้ผลการจำแนกประเภทเมื่อเทียบกับผลการแข่งขันจริง ดังตารางที่ \ref{tab:model_comparison}

\begin{table}[H]
\centering
\caption{เปรียบเทียบประสิทธิภาพของแบบจำลองในการทำนายผลการแข่งขัน (380 แมตช์)}
\label{tab:model_comparison}
\begin{tabular}{|l|c|c|c|c|c|}
\hline
\textbf{แบบจำลอง (Model)} & \textbf{Accuracy} & \textbf{Precision} & \textbf{Recall} & \textbf{F1-Score} & \textbf{Log Loss} \\ \hline
Baseline (Bookmaker Implied) & 52.63\% & 0.491 & 0.526 & 0.498 & 0.961 \\ \hline
Logistic Regression & 51.05\% & 0.472 & 0.511 & 0.478 & 0.985 \\ \hline
Random Forest & 53.16\% & 0.502 & 0.532 & 0.508 & 0.954 \\ \hline
XGBoost & 54.47\% & 0.518 & 0.545 & 0.523 & 0.942 \\ \hline
\textbf{LightGBM (Proposed)} & \textbf{55.79\%} & \textbf{0.531} & \textbf{0.558} & \textbf{0.537} & \textbf{0.928} \\ \hline
\end{tabular}
\end{table}

% -----------------------------------------------------------
\section{ผลการกระจายตัวการทำนายและ Confusion Matrix}
% -----------------------------------------------------------

เมื่อพิจารณาแบบจำลอง LightGBM ซึ่งให้ผลการทำนายแม่นยำที่สุด นำมาสร้าง Confusion Matrix แสดงรายละเอียดการจำแนกประเภทระหว่าง เจ้าบ้านชนะ (H), เสมอ (D), และ ทีมเยือนชนะ (A) ดังตารางที่ \ref{tab:confusion_matrix}

\begin{table}[H]
\centering
\caption{Confusion Matrix ของแบบจำลอง LightGBM สำหรับชุดข้อมูล Test Set (380 แมตช์)}
\label{tab:confusion_matrix}
\begin{tabular}{|c|c|c|c|c|}
\hline
\multicolumn{2}{|c|}{\multirow{2}{*}{\textbf{การทำนาย / ผลจริง}}} & \multicolumn{3}{c|}{\textbf{ผลการแข่งขันจริง (Actual)}} \\ \cline{3-5} 
\multicolumn{2}{|c|}{} & \textbf{เจ้าบ้านชนะ (H)} & \textbf{เสมอ (D)} & \textbf{ทีมเยือนชนะ (A)} \\ \hline
\multirow{3}{*}{\textbf{ผลทำนาย (Predicted)}} & \textbf{เจ้าบ้านชนะ (H)} & \textbf{138} & 28 & 19 \\ \cline{2-5} 
 & \textbf{เสมอ (D)} & 18 & \textbf{19} & 22 \\ \cline{2-5} 
 & \textbf{ทีมเยือนชนะ (A)} & 21 & 30 & \textbf{55} \\ \hline
\multicolumn{2}{|c|}{\textbf{รวมจำนวนแมตช์จริง}} & 177 & 77 & 126 \\ \hline
\end{tabular}
\end{table}

จากตารางที่ \ref{tab:confusion_matrix} สรุปผลการทำนายย่อยได้ดังนี้
\begin{enumerate}[label=\arabic*)]
    \item \textbf{การทำนายเจ้าบ้านชนะ (H):} ระบบทำนายได้ถูกต้อง 138 แมตช์ จากทั้งหมด 177 แมตช์ (Recall = 77.97\%) ซึ่งเป็นกลุ่มที่ทำนายได้แม่นยำที่สุด
    \item \textbf{การทำนายผลเสมอ (D):} ระบบทำนายได้ถูกต้องเพียง 19 แมตช์ จากทั้งหมด 77 แมตช์ (Recall = 24.68\%) เนื่องจากผลเสมอมักมีความผันผวนสูงเชิงสถิติ
    \item \textbf{การทำนายทีมเยือนชนะ (A):} ระบบทำนายได้ถูกต้อง 55 แมตช์ จากทั้งหมด 126 แมตช์ (Recall = 43.65\%)
\end{enumerate}

% -----------------------------------------------------------
\section{การวิเคราะห์ความสำคัญของตัวแปร (Feature Importance)}
% -----------------------------------------------------------

จากการวิเคราะห์ Feature Importance ของแบบจำลอง LightGBM พบว่าฟีเจอร์ที่มีผลต่อการทำนายมากที่สุด 3 อันดับแรก ได้แก่
\begin{enumerate}[label=\arabic*)]
    \item \textbf{elo\_diff:} ผลต่างคะแนน Elo Rating ระหว่างสองทีม (มีน้ำหนักการทำนาย 28.4\%)
    \item \textbf{xg\_moving\_avg\_5:} ค่าเฉลี่ยประตูที่คาดว่าจะเกิดขึ้น (xG) ย้อนหลัง 5 นัด (มีน้ำหนัก 21.2\%)
    \item \textbf{key\_injured\_count:} จำนวนผู้เล่นตัวหลักที่บาดเจ็บ (มีน้ำหนัก 14.5\%)
\end{enumerate}

% -----------------------------------------------------------
\section{การวิเคราะห์ผลการทดลองและข้อผิดพลาด}
% -----------------------------------------------------------

จากผลการทดลองแสดงให้เห็นว่าแบบจำลอง LightGBM ให้ค่า Accuracy สูงถึง 55.79\% ซึ่งสูงกว่า Baseline ของราคาต่อรองเจ้ามือ (52.63\%) อย่างไรก็ตาม การวิเคราะห์ข้อผิดพลาด (Error Analysis) พบว่าข้อจำกัดสำคัญของระบบคือการทำนายผลเสมอ (Draw) เนื่องจากแมตช์ที่จบลงด้วยการเสมอมักเกิดจากปัจจัยที่ไม่คาดคิด เช่น ใบแดงระหว่างเกม หรือความผิดพลาดส่วนบุคคลของผู้เล่น ซึ่งข้อมูลสถิติล่วงหน้าไม่สามารถครอบคลุมได้ทั้งหมด